\documentclass[10pt,a4paper]{amsart}
\usepackage[margin=19mm]{geometry}
\usepackage{amsmath,amssymb,mathtools}
\usepackage[scr=rsfs,cal=euler]{mathalfa}
\usepackage{xfrac}
\usepackage[unicode,hidelinks,bookmarks=false]{hyperref}
\newcommand{\ie}{i.e.\ }
\newcommand{\cf}{cf.\ }
\newcommand{\resp}{respectively\ }
\newcommand{\Subsection}{\subsection}
\providecommand{\nobreakdash}{\nobreak\hbox{-}}
\setlength{\emergencystretch}{2em}
\multlinegap0pt
\newcommand*\samethanks[1][\value{footnote}]{\footnotemark[#1]}
\usepackage{mathtools}
\usepackage{dsfont}
\usepackage{float}
\usepackage{bm}
\usepackage{amssymb}
\usepackage{stackengine}
\newtheorem{theorem}{Theorem}
\newcommand{\OA}{\mathcal{M}_{\mathrm{OA}}}
\newcommand{\R}{\mathbb{R}}
\newcommand{\T}{\mathbb{T}}
\newcommand{\Z}{\mathbb{Z}}
\renewcommand{\d}{\textnormal{d}}
\newcommand{\jump}{\hfill\break}
\renewcommand{\u}{\textnormal{u}}
\title{Unstable Manifolds for the Kuramoto Model:\\ Convergence to the Ott-Antonsen Manifold}
\author{Christian Kuehn, Giacomo Landi\samethanks}
\date{Source: arXiv:2608.24453v2 (CC BY 4.0)}
\begin{document}
\maketitle

\noindent\emph{Source and licence.} Unstable Manifolds for the Kuramoto Model:\\ Convergence to the Ott-Antonsen Manifold. Version arXiv:2608.24453v2 (CC BY 4.0), distributed by arXiv under the Creative Commons Attribution 4.0 International licence, http://creativecommons.org/licenses/by/4.0/, as recorded in the arXiv licence metadata for this exact version and shown on the abstract page https://arxiv.org/abs/2608.24453v2. Full licence notice: \texttt{licenses/arxiv-2608.24453-CC-BY-4.0.txt}.\par\smallskip

\noindent\textit{Editorial scope.} Counted unit: the article's theorem conv\_rate\_synchro (source0.tex lines 745--756). The article's complete original proof of it is delivered with it (source0.tex lines 757--866, one \texttt{proof} environment of 110 lines). No mathematics is written, repaired or re-derived: the statement and the proof are the article's own lines, in the article's own notation, and no retained line is substituted or re-flowed.  Retained uncounted as the source-local context the proof needs: lines 149--151; lines 489--500; lines 493--499; lines 539--541; lines 665--671.  Nothing was omitted from any retained span, so the package carries no omission record.  No retained line contains a citation, so the wrapper carries no bibliography and no bibliography entry.  No retained line contains a figure environment, an image-inclusion command or drawing code, so no image is delivered and the package declares no asset dependency.  Editorial formatting only, in addition: an AMS \texttt{amsart} preamble that restates the article's own declarations and macros used by the retained lines, namely the environments \texttt{theorem} on separate counters, together with the standard packages those lines need.  The retained environments print in this document as Theorem 1 in order of appearance, whereas the article prints the same items with the numbering of its own full document.  The complete native source of the article as distributed in the arXiv e-print for this exact version is bundled unchanged as \texttt{sources/208478/source0.tex}.
\begin{equation}\label{KMH}
    \dot{\theta}_i(t)=\frac{K}{N}\sum_{j=1}^N\sin{(\theta_j(t)-\theta_i(t))}.\qquad\forall i=1,...,N\tag{KMH}
\end{equation}

\begin{theorem}\label{main_theorem}
    Every element of the family of incoherent equilibria $\theta^N_q$ of \eqref{KMH}, with $q\in\mathbb{T}$ and $N\geq3$, admits an explicit expression for its two-dimensional unstable manifold, namely
    $$W^\u(\theta^N_q)=\left\{\theta\in\T^N\text{ s.t. } \theta_i=\phi_i(\alpha,\beta,A);\text{ }\frac{1}{N}\sum_{i=1}^N\theta_i=q\text{ }\bigg\vert\;\beta\in[0,1),\;\alpha\in[-\pi,\pi],A\in\left[-\frac{\pi}{2},\frac{\pi}{2}\right]\right\}.$$
    Moreover, $(\beta,\alpha,A)$ satisfy the following system of ODEs:
    \begin{equation}\label{3d-system}
    \begin{cases}
        \dot{\beta}=\frac{K}{2}\beta(1-\beta^2)B_N(A,\beta),\\
        \dot{\alpha}=\frac{K}{2}(1+\beta^2)D_N(A,\beta),\\
        \dot{A}=\frac{K}{4}(1-\beta^2)D_N(A,\beta).
    \end{cases}
    \end{equation}
\end{theorem}

    \begin{equation}
    \begin{cases}
        \dot{\beta}=\frac{K}{2}\beta(1-\beta^2)B_N(A,\beta),\\
        \dot{\alpha}=\frac{K}{2}(1+\beta^2)D_N(A,\beta),\\
        \dot{A}=\frac{K}{4}(1-\beta^2)D_N(A,\beta).
    \end{cases}
    \end{equation}

\begin{equation}\label{mean_phase_new}
    -\alpha+\frac{2}{N}\sum_{i=1}^N\arctan\left(\frac{1-\beta}{1+\beta}\tan\left(\pi\xi_i+A\right)\right)=q.
\end{equation}

\begin{equation}\label{limit_dyn}
    \begin{cases}
    \dot{\beta}=\frac{K}{2}\beta(1-\beta^2),\\
    \dot{\alpha}=0,\\
    \dot{A}=0,
\end{cases}
\end{equation}

\begin{theorem}\label{conv_rate_synchro}
\begin{enumerate} 
\item[i)] For every $(\alpha(0),\beta(0))\in[-\pi,\pi]\times(0,1)$, the trajectory $f_{\alpha(t),\beta(t)}\in\OA$ converges exponentially towards $\delta_{\alpha(0)}$ with respect to the $p$-Wasserstein distance. More precisely, for every $p\geq1$ there exists a constant $\Tilde{C}_p>0$ such that, for every initial datum $(\alpha(0),\beta(0))\in[-\pi,\pi]\times(0,1)$ and every $t\in\R^+$, the following estimates hold:
\begin{equation} W_p\big(f_{\alpha(t),\beta(t)}\d\theta,\delta_{\alpha(0)}\big)^p \leq \Tilde{C}_p\frac{1-\beta(0)^2}{\beta(0)^2} e^{-Kt}, \qquad \forall p>1; \end{equation} 
\begin{equation} W_1\big(f_{\alpha(t),\beta(t)}\d\theta,\delta_{\alpha(0)}\big) \leq \Tilde{C}_1 \frac{1-\beta(0)^2}{\beta(0)^2} e^{-Kt} \left( 1+\left|\ln\frac{1-\beta(0)^2}{\beta(0)^2}\right|+Kt \right). 
\end{equation} 
\item[ii)] For every $q\in\T^1$, $N\geq3$ and $\theta(0)\in W^\u(\theta^N_q)$, one has \[ \lim_{t\rightarrow+\infty}|\theta(t)-q\underline{1}|=0. \] Moreover, for every $p\geq1$ there exists a constant $\Tilde{D}_p>0$ such that, for every $t\in\R$, the following estimates hold:
\begin{equation} W_p\left(\mu^N(\theta(t)),\delta_{q}\right)^p \leq \Tilde{D}_p \left( \frac{1}{N} + \frac{1-\beta(0)^2}{\beta(0)^2} e^{-\frac{K}{3}t} \right), \qquad \forall p>1; \end{equation}
\begin{equation} W_1\left(\mu^N(\theta(t)),\delta_{q}\right) \leq \Tilde{D}_1 \left( \frac{1}{N} + \frac{1-\beta(0)^2}{\beta(0)^2} e^{-\frac{K}{3}t} \left( 1+\left|\ln \frac{1-\beta(0)^2}{\beta(0)^2}\right| +\frac{K}{3}t \right) \right).
\end{equation}
\end{enumerate}
\end{theorem}

    \begin{proof} \textbf{i)} By the translational invariance of the Wasserstein distance $W_p$ induced by the Euclidean metric, we may assume without loss of generality that $\alpha=0$. Since the product measure is an admissible coupling between $f_{0,\beta}\d\theta$ and $\delta_0$, we obtain 
    $$W_p(f_{0,\beta}\d\theta,\delta_0)^p\leq\int_{-\pi}^\pi|\theta|^pf_{0,\beta}(\theta)\d\theta =2\int_0^\pi\theta^p\frac{1-\beta^2}{2\pi(1-2\beta\cos\theta+\beta^2)}\,\d\theta.$$
    For $\beta\in[0,1/2]$, it follows immediately that
    \begin{equation}\label{low_beta} W_p(f_{0,\beta}\d\theta,\delta_0)^p\leq\pi^p.
    \end{equation}
    On the other hand, for $\beta\in[1/2,1)$ we have
    \begin{align*}
            W_p(f_{0,\beta}\d\theta,\delta_0)^p&\leq\frac{1}{\pi}\int_0^\pi\theta^p\frac{1-\beta^2}{1-2\beta\cos\theta+\beta^2}\d\theta=\frac{1}{\pi}\int_0^\pi\theta^p\frac{1-\beta^2}{(1-\beta)^2+2\beta(1-\cos\theta)}\d\theta\leq\\
            &\leq\frac{1}{\pi}\int_0^\pi\theta^p\frac{1-\beta^2}{(1-\beta)^2+\frac{2}{\pi^2}\theta^2}\d\theta\leq\frac{2}{\pi}(1-\beta)\int_0^\pi\frac{\theta^p}{(1-\beta)^2+\frac{2}{\pi^2}\theta^2}\d\theta\leq\\
            &\leq\pi(1-\beta)\int_0^\pi\frac{\theta^p}{(1-\beta)^2+\theta^2}\d\theta
    \end{align*}
    where we used the inequality $$1-\cos\theta\geq\frac{2}{\pi^2}\theta^2,\qquad \forall\theta\in[-\pi,\pi].$$ We now split the last integral as 
    $$I_1\coloneqq\int_0^{1-\beta}\frac{\theta^p}{(1-\beta)^2+\theta^2}\,\d\theta, \qquad I_2\coloneqq\int_{1-\beta}^{\pi}\frac{\theta^p}{(1-\beta)^2+\theta^2}\,\d\theta.$$
    First, for every $p\geq1$,
    $$I_1\leq\int_0^{1-\beta}\frac{\theta^p}{(1-\beta)^2}\,\d\theta =\frac{(1-\beta)^{p-1}}{p+1}.$$
    Moreover, for $p>1$,
    $$I_2\leq\int_{1-\beta}^{\pi}\theta^{p-2}\,\d\theta =\frac{\pi^{p-1}-(1-\beta)^{p-1}}{p-1} \leq\frac{\pi^{p-1}-2^{1-p}}{p-1},$$
    while for $p=1$,
    $$I_2\leq|\ln(1-\beta)|+\ln\pi.$$
    Therefore, for every $\beta\in[1/2,1)$ and every $p>1$,
    $$W_p(f_{0,\beta}\d\theta,\delta_0)^p \leq\hat{C}_p(1-\beta),$$
    where
    $$\hat{C}_p\coloneqq \pi\left(\frac{1}{p+1} + \frac{\pi^{p-1}-2^{1-p}}{p-1}\right).$$
    For $p=1$, we similarly obtain
    $$W_1(f_{0,\beta}\d\theta,\delta_0) \leq \Tilde{C}_1(1-\beta)\bigl(1+|\ln(1-\beta)|\bigr),$$
    where $$\Tilde{C}_1\coloneqq \pi\left(\frac12+\ln\pi\right).$$
    It remains to combine these estimates with the case $\beta\in[0,1/2]$. Since
    $$\frac12\leq1-\beta, \qquad \frac{1+\ln2}{2}\leq(1-\beta)\bigl(1+|\ln(1-\beta)|\bigr),$$
    estimate \eqref{low_beta} implies that, for $p>1$,
    $$W_p(f_{0,\beta}\d\theta,\delta_0)^p \leq2\pi^p(1-\beta),$$
    while for $p=1$,
    $$W_1(f_{0,\beta}\d\theta,\delta_0) \leq \frac{2\pi}{1+\ln2} (1-\beta)\bigl(1+|\ln(1-\beta)|\bigr).$$
    Combining the previous estimates, we conclude that for every $\beta\in[0,1)$ and every $p>1$,
    $$W_p(f_{0,\beta}\d\theta,\delta_0)^p \leq \Tilde{C}_p(1-\beta),$$
    where
    $$\Tilde{C}_p\coloneqq\max\{2\pi^p,\hat{C}_p\}.$$
    For $p=1$, we have
    $$W_1(f_{0,\beta}\d\theta,\delta_0) \leq \Tilde{C}_1(1-\beta)\bigl(1+|\ln(1-\beta)|\bigr).$$
    It remains to derive the exponential decay of $1-\beta(t)$. From \eqref{limit_dyn}, the explicit expression for $\beta(t)$ is
    $$\beta(t)=\frac{\beta(0)} {\sqrt{\beta(0)^2+(1-\beta(0)^2)e^{-Kt}}},$$
    which implies
    $$1-\beta(t) \leq 1-\beta(t)^2 = \frac{(1-\beta(0)^2)e^{-Kt}} {\beta(0)^2+(1-\beta(0)^2)e^{-Kt}} \leq \frac{1-\beta(0)^2}{\beta(0)^2}e^{-Kt}.$$
    Combining this estimate with the monotonicity of $h_1(x)=x$ and $h_2(x)=x(1+|\ln x|)$ completes the proof of the first statement. \jump\jump
    \textbf{ii)} As before we begin by estimating the Wasserstein distance between a generic point of $W^\u(\theta_q^N)$ and the synchronized state $\delta_q$ in terms of $\beta$. We then exploit the dynamics of $\beta_N(t)$ to turn this static estimate into a quantitative convergence result.\jump\jump\noindent
    \fbox{\textbf{Step 1:} Estimate for $W_p\left(\mu^N(\theta(t)),\delta_q\right)$} We fix $A$ and $\beta$ and, for brevity, set $y_i\coloneqq\pi\xi_i+A.$ Moreover, as in Step 1 of the proof of Theorem \ref{main_theorem}, we define
    $$\varphi_i(\beta,A)\coloneqq\phi_i(\alpha,\beta,A)+\alpha.$$
    Then, using again the first integral \eqref{mean_phase_new} and Jensen's inequality, we obtain 
    \begin{align*}
           W_p\left(\mu^N(\theta(t),\delta_{q})\right)^p&\leq\frac{1}{N}\sum_{i=1}^N\Big|\phi_i(\alpha_N(\beta,A),\beta,A)-q\Big|^p=\frac{1}{N}\sum_{i=1}^N\Big|\varphi_i(\beta,A)-\frac{1}{N}\sum_{j=1}^N\varphi_j(\beta,A)\Big|^p\leq\\
           &\leq\frac{2^{p-1}}{N}\sum_{i=1}^N\left(\big|\varphi_i(\beta,A)\big|^p+\frac{1}{N}\sum_{j=1}^N\big|\varphi_j(\beta,A)\big|^p\right)=2^p\frac{1}{N}\sum_{i=1}^N\big|\varphi_i(\beta,A)\big|^p
       \end{align*}
       Let
       $$r(y)\coloneqq d\left(y,\frac{\pi}{2}+\pi\Z\right),\qquad r_i\coloneqq r(y_i),\qquad \rho\coloneqq\frac{1-\beta}{1+\beta}.$$
       We claim that
       $$|\varphi_i|\leq\pi\min\left\{1,\frac{\rho}{r_i}\right\}.$$
       Indeed, by definition one directly has $|\varphi_i|\leq\pi$. On the other hand, using the Lipschitz continuity of $\arctan(\cdot)$, we get
       $$ |\varphi_i| \leq 2\rho|\tan(y_i)| = 2\rho\frac{|\sin(y_i)|}{|\cos(y_i)|} \leq \frac{2\rho}{|\cos(y_i)|} = \frac{2\rho}{\sin(r_i)} \leq \pi\frac{\rho}{r_i}, $$
       where we used that $\sin x\geq\frac{2}{\pi}x$ for every $x\in[0,\pi/2]$. Therefore,
       $$W_p\left(\mu^N(\theta(t)),\delta_q\right)^p \leq \frac{2^p\pi^p}{N} \sum_{i=1}^N \min\left\{1,\frac{\rho^p}{r_i^p}\right\}.$$
       We now derive a lower bound for the quantities $r_i$. Since the points $y_i$ form a uniform grid on $\R/\pi\Z$, with spacing $\pi/N$, we have
       $$ \#\left\{i:r_i\leq R\right\} \leq \frac{2R}{\pi/N}+1. $$
       Up to relabeling, we may assume without loss of generality that
       $$0\leq r_1\leq r_2\leq\cdots\leq r_N.$$
       Thus, for every $m=1,\ldots,N$,
       $$ \#\left\{i:r_i\leq r_m\right\}\geq m. $$
       Combining the last two inequalities gives
       $$r_m\geq\pi\frac{m-1}{2N}.$$
       Using this bound for all terms in the previous summation except the first one, we obtain
       \begin{align*}
        W_p\left(\mu^N(\theta(t),\delta_{q})^p\right)&\leq2^p\pi^p\left(\frac{1}{N}+\frac{1}{N}\sum_{n=2}^N\min\left\{1,\left(\frac{a}{n-1}\right)^p\right\}\right)\leq\\
        &\leq2^p\pi^p\left(\frac{1}{N}+\frac{1}{N}\sum_{n=1}^N\min\left\{1,\left(\frac{a}{n}\right)^p\right\}\right)
       \end{align*}
       with $a\coloneqq\frac{2N\rho}{\pi}$. If $a\in[0,1)$ and $p>1$, then
       $$ \sum_{n=1}^N \min\left\{1,\left(\frac{a}{n}\right)^p\right\} = \sum_{n=1}^N\left(\frac{a}{n}\right)^p \leq a^p\zeta(p) \leq a\zeta(p), $$
       where $\zeta(\cdot)$ denotes the Riemann zeta function. For $p=1$, we use
       $$\sum_{n=1}^N\frac{1}{n}\leq1+\ln N,$$
       and hence
       $$ \sum_{n=1}^N \min\left\{1,\frac{a}{n}\right\} \leq a(1+\ln N). $$
       On the other hand, if $a\geq1$ and $p>1$, then 
       \begin{align*}
            \sum_{n=1}^N\min\left\{1,\left(\frac{a}{n}\right)^p\right\}&=\sum_{n=1}^{\lfloor a\rfloor}\min\left\{1,\left(\frac{a}{n}\right)^p\right\}+\sum_{n=\lfloor a\rfloor+1}^N\min\left\{1,\left(\frac{a}{n}\right)^p\right\}=\\
           &=\lfloor a\rfloor +a^p\sum_{n=\lfloor a\rfloor+1}^N n^{-p}\leq a+a^p\int_{\lfloor a\rfloor}^\infty x^{-p}\d x=\\
           &=a+a^p\frac{\lfloor a \rfloor^{1-p}}{p-1}\leq a+a^p\frac{a^{1-p}2^{p-1}}{p-1}= a\left(1+\frac{2^{p-1}}{p-1}\right).
        \end{align*}
        Similarly, for $p=1$ we obtain
        \begin{align*}
            \sum_{n=1}^N\min\left\{1,\frac{a}{n}\right\}&=\sum_{n=1}^{\lfloor a\rfloor}\min\left\{1,\frac{a}{n}\right\}+\sum_{n=\lfloor a\rfloor+1}^N\min\left\{1,\frac{a}{n}\right\}=\\
            &=\lfloor a\rfloor +a\sum_{n=\lfloor a\rfloor+1}^N \frac{1}{n}\leq a+a\int_{\lfloor a\rfloor}^N \frac{1}{x}\d x=a\left(1+\ln\frac{N}{\lfloor a\rfloor}\right)\leq a\left(1+\ln\frac{2N}{a}\right).
        \end{align*}
        Combining the previous estimates, for $p>1$ we have
        $$ \sum_{n=1}^N \min\left\{1,\left(\frac{a}{n}\right)^p\right\} \leq a\max\left\{\zeta(p),1+\frac{2^{p-1}}{p-1}\right\}. $$
        Since $\rho\leq1-\beta$, this gives 
        $$ \sum_{n=1}^N \min\left\{1,\left(\frac{a}{n}\right)^p\right\} \leq (1-\beta)\frac{2N}{\pi} \max\left\{\zeta(p),1+\frac{2^{p-1}}{p-1}\right\}. $$
        Therefore,
        $$ W_p\left(\mu^N(\theta(t)),\delta_q\right)^p \leq \Tilde{D}_p \left( \frac{1}{N} + 1-\beta \right), $$
        where
        $$ \Tilde{D}_p\coloneqq 2^{1+p}\pi^{p-1} \max\left\{\zeta(p),1+\frac{2^{p-1}}{p-1}\right\}. $$
        In the case $p=1$, the same argument yields
        $$ \sum_{n=1}^N \min\left\{1,\frac{a}{n}\right\} \leq \frac{2N}{\pi} \left(1+\ln(2\pi)\right) (1-\beta) \left( 1+\ln\left(\frac{1}{1-\beta}\right) \right), $$
        and consequently
        $$ W_1\left(\mu^N(\theta(t)),\delta_q\right) \leq \Tilde{D}_1 \left( \frac{1}{N} + (1-\beta) \left( 1+\ln\left(\frac{1}{1-\beta}\right) \right) \right), $$
        with
        $$ \Tilde{D}_1\coloneqq4(1+\ln(2\pi)). $$\jump\noindent \fbox{\textbf{Step 2:} Exponential decay for $\beta_N(t)$} It remains to prove that $\beta_N(t)$ converges exponentially to $1$. From the first equation of \eqref{3d-system} and the estimate $B_N(A,\beta)\geq\frac{N-2}{N}$, obtained in Step 2 of Theorem \ref{main_theorem}, we get $$ \dot{\beta}_N \geq \frac{K}{6}\beta_N(1-\beta_N^2) $$ for all $N\geq3$. Let $x(t)$ be the solution of
        $$ \dot{x} = \frac{K}{6}x(1-x^2), \qquad x(0)=\beta_N(0). $$
        Then
        $$ 1-x(t) \leq \frac{1-x(0)^2}{x(0)^2} e^{-\frac{K}{3}t}. $$
        By the comparison principle, since $\beta_N(0)=x(0)$, we have $\beta_N(t)\geq x(t)$ for all $t\geq0$. Hence,
        $$ 1-\beta_N(t) \leq 1-x(t) \leq \frac{1-\beta_N(0)^2}{\beta_N(0)^2} e^{-\frac{K}{3}t}. $$
        Combining this decay estimate with the bounds obtained in Step 1 proves the second statement of the theorem.
        \end{proof}

\end{document}
